\chapter{Bài toán}
\label{ch:problem}

\section{Vì sao mỗi tác vụ một mô hình riêng khó mở rộng}

Các mô hình ngôn ngữ huấn luyện trước (pre-trained) đã thay đổi cách xây dựng hệ thống xử lý ngôn ngữ tự nhiên. Một encoder như BERT~\cite{devlin2019bert} có thể được tinh chỉnh (fine-tune) cho một bài toán phân loại mới với lượng dữ liệu hạn chế, nên một hệ thống hỗ trợ nhiều tác vụ hiểu ngôn ngữ thường tích lũy một bản sao đã tinh chỉnh của cùng một kiến trúc cho mỗi tác vụ. Mỗi bản sao hoạt động tốt trên tác vụ của nó, nhưng cả tập hợp trở nên tốn kém theo bốn cách.

\begin{itemize}
  \item \textbf{Lưu trữ.} Mỗi bản sao giữ một bộ tham số hoàn chỉnh, dù tất cả đều xuất phát từ cùng một bộ trọng số huấn luyện trước.
  \item \textbf{Huấn luyện.} Mỗi tác vụ mới cần thêm một lần tinh chỉnh đầy đủ, với dữ liệu, tài nguyên tính toán và quy trình chọn mô hình riêng.
  \item \textbf{Vận hành.} Hệ thống phải quản lý phiên bản, nạp và chuyển yêu cầu đến một mô hình riêng cho từng tác vụ.
  \item \textbf{Cô lập kiến thức.} Mỗi mô hình chỉ biết tác vụ của nó. Những gì mô hình học được khi tinh chỉnh không đến được các mô hình khác.
\end{itemize}

\section{Phát biểu bài toán và câu hỏi nghiên cứu}
\label{sec:problem}

Gọi $\theta_0$ là trọng số của một mô hình huấn luyện trước, và $\theta_i$, $i\in\mathcal{T}$, là trọng số của các mô hình được tinh chỉnh từ nó cho một tập $\mathcal{T}$ gồm $T$ tác vụ. Chúng tôi gọi mỗi mô hình đã tinh chỉnh là một \emph{chuyên gia} (expert). \emph{Véc-tơ tác vụ} (task vector) của tác vụ $i$ là phần thay đổi mà quá trình tinh chỉnh tạo ra so với điểm xuất phát chung:
\begin{equation}
\tau_i=\theta_i-\theta_0 .
\label{eq:task-vector}
\end{equation}
Đầu vào của bài toán là $\theta_0$ cùng các trọng số $\theta_i$, hoặc tương đương là các véc-tơ tác vụ. Đầu ra là một mô hình đa nhiệm với trọng số $\theta_{\mathrm{mer}}=C(\theta_0,\theta_i)$, ánh xạ một văn bản và tác vụ của nó thành một nhãn. Bài toán là tìm hàm hợp nhất $C$ sao cho mô hình hợp nhất đạt ba mục tiêu:

\begin{enumerate}
  \item[\textbf{O1.}] Tích hợp kiến thức của các tác vụ một cách \emph{động}, tùy theo đầu vào.
  \item[\textbf{O2.}] Giữ chất lượng gần với các mô hình riêng theo tác vụ.
  \item[\textbf{O3.}] Cần ít tham số lưu trữ hơn so với việc giữ tất cả các mô hình.
\end{enumerate}

Không có ngưỡng đạt nào được đặt trước, nên ``gần'' được báo cáo như một khoảng cách đo được chứ không phải một kết quả đúng hoặc sai. Nghiên cứu này giải quyết bốn câu hỏi:

\begin{enumerate}
  \item[\textbf{RQ1.}] Các bộ định tuyến có thể khôi phục bao nhiêu hiệu năng bị mất do hợp nhất tĩnh, khi được huấn luyện trên các mẫu có nhãn mà các chuyên gia chưa từng thấy?
  \item[\textbf{RQ2.}] Mức khôi phục đó phụ thuộc vào ngân sách mẫu có nhãn như thế nào?
  \item[\textbf{RQ3.}] Cải thiện này có đòi hỏi định tuyến phụ thuộc đầu vào không, hay chỉ cần các hệ số cố định tốt hơn?
  \item[\textbf{RQ4.}] Việc chọn các lớp MLP để nâng cấp ảnh hưởng thế nào đến hành vi theo tác vụ và dung lượng lưu trữ, và phương pháp tốn bao nhiêu khi suy luận?
\end{enumerate}

\section{Vì sao bài toán quan trọng và vì sao khó}

\emph{Hợp nhất mô hình} (model consolidation), tức gộp các mô hình đã tinh chỉnh thành một, quan trọng vì nhiều hệ thống phục vụ các tác vụ ngôn ngữ có liên quan với nhau, và một mô hình chung giúp tái sử dụng kiến thức giữa chúng. Học đa nhiệm (multi-task learning)~\cite{caruana1997multitask} cũng tạo ra một mô hình, nhưng phải huấn luyện đồng thời trên toàn bộ dữ liệu và lặp lại chi phí đó khi thêm tác vụ. Hợp nhất mô hình bắt đầu từ các mô hình đã có, nên không cần huấn luyện chung lại từ đầu.

Cái khó nằm ở cách các mô hình đã tinh chỉnh khác biệt với nhau. Cập nhật có lợi cho tác vụ này có thể gây hại cho tác vụ khác khi cùng tác động lên một tham số, gọi là \emph{can nhiễu giữa các tác vụ}. Hai mô hình cũng có thể kéo một trọng số theo hai hướng ngược nhau, nên phép cộng làm mất một phần những gì mỗi mô hình đã học. Vì vậy, mọi phép hợp nhất đều phải quyết định giữ lại bao nhiêu từ mỗi cập nhật, và phải làm điều đó mà không có dữ liệu huấn luyện gốc của mọi tác vụ.

\section{Vì sao chỉ hợp nhất tĩnh có thể chưa đủ}

Một phương pháp tĩnh như Task Arithmetic~\cite{ilharco2023task} kết hợp tham số một lần và dùng một bộ trọng số cho mọi đầu vào. Năm tác vụ trong nghiên cứu này làm rõ giới hạn đó. Chúng đòi hỏi những điều khác nhau ở encoder: đánh giá ngữ pháp hoặc cảm xúc của một câu, hoặc xét quan hệ giữa hai câu qua diễn đạt lại, suy diễn hay trả lời câu hỏi. Một véc-tơ tham số duy nhất phải phục vụ tất cả cùng lúc, và hệ số tỉ lệ của nó chỉ có thể đánh đổi chúng với nhau. Như Tang và cộng sự~\cite{tang2024wemoe} lập luận, nghiệm tốt nhất có thể thay đổi theo đầu vào, nên một kết quả cố định khó có thể tốt nhất cho mọi tác vụ. Vì vậy, khoảng trống mà báo cáo này xem xét là liệu trọng số hợp nhất có nên thay đổi theo đầu vào thay vì cố định hay không.
